% =====================================================================
%  PAPER II -- broadened from proposal P3 (certificate as compliance
%  artifact). Target: IEEE S&P, Systematization of Knowledge (SoK) track.
%
%  Standalone. Set as Overleaf's main document; pdfLaTeX + BibTeX.
%    pdflatex PaperII_sok_quantifier_gap && bibtex PaperII_sok_quantifier_gap \
%      && pdflatex PaperII_sok_quantifier_gap && pdflatex PaperII_sok_quantifier_gap
% =====================================================================
\documentclass[11pt,a4paper]{article}
% =====================================================================
%  Shared preamble for the three research proposals.
%
%  Each proposal \input{}s this file and is otherwise standalone, so any
%  one of them can be set as Overleaf's main document and compiled on its
%  own. Nothing here is specific to a single proposal.
% =====================================================================
\usepackage[T1]{fontenc}
\usepackage[utf8]{inputenc}
% Times, matching the parent manuscript. mathptmx is in every TeX Live and on
% Overleaf; lmodern is deliberately not used, since it is absent from minimal
% distributions and this preamble should compile anywhere.
\usepackage{mathptmx}
\usepackage[a4paper,margin=25mm]{geometry}
\usepackage{amsmath,amssymb,amsthm}
\usepackage{booktabs}
\usepackage{tabularx}
\usepackage{xcolor}
\usepackage{titlesec}
\usepackage{enumitem}
\usepackage{microtype}
\usepackage[hidelinks,breaklinks]{hyperref}

% ---- shared style tokens, matching the parent manuscript ------------
\definecolor{netcol}{RGB}{31,78,121}
\definecolor{autocol}{RGB}{18,120,90}
\definecolor{bridgecol}{RGB}{176,96,0}
\definecolor{accent}{RGB}{140,20,20}

\newcommand{\layernet}{\textcolor{netcol}{\textbf{network}}}
\newcommand{\layerauto}{\textcolor{autocol}{\textbf{autonomy}}}
\newcommand{\dataset}[1]{\textsc{#1}}
\newcommand{\code}[1]{\texttt{#1}}
\newcommand{\Real}{\mathbb{R}}
\newcommand{\Norm}[1]{\left\lVert #1\right\rVert}
\newcommand{\MCR}{\mathrm{MCR}}

\newtheorem{claim}{Claim}
\newtheorem{proposition}{Proposition}
\newtheorem{definition}{Definition}
\theoremstyle{remark}
\newtheorem{remark}{Remark}

% ---- section styling -------------------------------------------------
\titleformat{\section}{\large\bfseries\color{netcol}}{\thesection}{0.7em}{}
\titleformat{\subsection}{\normalsize\bfseries}{\thesubsection}{0.6em}{}
\titlespacing*{\section}{0pt}{1.4em}{0.6em}

\setlist[itemize]{leftmargin=1.35em,topsep=0.3em,itemsep=0.25em}
\setlist[enumerate]{leftmargin=1.6em,topsep=0.3em,itemsep=0.25em}

% ---- a boxed "what this rests on" callout ---------------------------
\newenvironment{honestly}
  {\par\medskip\noindent\begin{tabular}{@{}p{2.5mm}@{\hspace{3mm}}p{0.9\linewidth}@{}}
   \textcolor{accent}{\rule{2.5mm}{\baselineskip}} &
   \small\itshape}
  {\end{tabular}\par\medskip}

% ---- the proposal header --------------------------------------------
\newcommand{\proposalheader}[3]{%
  \begin{center}
  {\Large\bfseries #1\par}\medskip
  {\large\itshape #2\par}\medskip
  {\small #3\par}
  \end{center}
  \medskip\hrule\medskip}

\usepackage{array}
\newcolumntype{Y}{>{\raggedright\arraybackslash}X}
\newcolumntype{P}[1]{>{\raggedright\arraybackslash}p{#1}}
% Quantifier-level labels. Serif rather than \textsf: the sans face is a
% bitmap font in minimal TeX installs, and microtype cannot expand it.
\newcommand{\Qlvl}[1]{\textbf{Q}\textsubscript{\textit{#1}}}
\newcommand{\Qin}{\Qlvl{in}}
\newcommand{\Qtr}{\Qlvl{tr}}
\newcommand{\Qop}{\Qlvl{op}}
\newcommand{\Qpop}{\Qlvl{pop}}

\begin{document}

\proposalheader
  {Paper II --- SoK}
  {SoK: The Quantifier Gap Between Certified Robustness and Airworthiness
   Assurance}
  {Integrated proposal, broadened from P3 $\cdot$ RobustUAVs.ai $\cdot$
   Roger Nick Anaedevha $\cdot$ 2026-10-01\\
   Target: IEEE Symposium on Security and Privacy, SoK track $\cdot$
   fallback: DASC or ACM Computing Surveys}

\begin{abstract}
\noindent
Certified-robustness research and aviation assurance both promise guarantees
about learned systems, and both are now being asked to meet: EASA's
learning-assurance framework and the forthcoming ED-324 expect evidence about
machine-learning constituents, and the robustness literature produces
certificates in quantity. We argue that the two rarely meet because they
quantify over different things. A certified radius is a statement about
\emph{every perturbation of one input}; a contingency volume is a statement
about \emph{one operation}; a safety objective per flight hour is a statement
about \emph{a population of operations}. Evidence transfers from a guarantee to
an objective only when the guarantee's quantifier dominates the objective's, or
when a \emph{lifting argument} --- a dynamics bound, a concentration inequality,
a runtime monitor --- bridges them, at the price of an assumption that itself
needs evidence. We propose to systematise the certified-guarantee literature for
learned perception, navigation and control by quantifier, to code SORA and EASA
learning-assurance objectives the same way, and to tabulate which transfers are
sound, which require a lift, and which are category errors. A SORA Annex~A
evidence generator demonstrates the sound path end to end.
\end{abstract}

% =====================================================================
\section{The observation}
% =====================================================================
P3 proposed mapping one certificate --- the parent manuscript's~\cite{parentpaper}
--- onto SORA and ED-324 vocabulary, and its own honesty note predicted the
likely finding: \emph{certified radii are per-input, while assurance objectives
are per-operation}. That sentence is more general than P3 used it for. It is a
claim about an entire literature, and it is testable.

It is also incomplete in one direction, which is what makes it worth a
systematisation rather than a remark. Aviation does not quantify only per
operation. A SORA contingency volume is sized for one operation; but a safety
objective expressed per flight hour, and EASA's learning-assurance expectation
that a model generalises across its Operational Design Domain, are statements
about \emph{populations}. So the two communities do not simply sit at two ends of
one axis. Each produces statements at several quantifier levels, and the
question is which pairs line up.

% =====================================================================
\section{The framework: four quantifier levels}
\label{sec:framework}
% =====================================================================
\begin{center}\small
\begin{tabularx}{\linewidth}{@{}P{0.12\linewidth}P{0.25\linewidth}YY@{}}
\toprule
\textbf{Level} & \textbf{Quantifies over} & \textbf{Typical guarantee} &
\textbf{Typical assurance objective} \\
\midrule
\Qin  & all perturbations of one input &
certified radius: randomized smoothing~\cite{cohen2019smoothing}, bound
propagation~\cite{wang2021betacrown}, SMT~\cite{katz2017reluplex} &
rarely stated at this level \\
\Qtr  & all disturbances along one trajectory, over a horizon &
reachability tube; Lipschitz--Gr\"onwall~\cite{gronwall1919} &
position-holding error; flight technical error \\
\Qop  & one complete operation &
mission predicate (spatial \emph{and} temporal); runtime assurance with
handover~\cite{sha2001simplex} &
SORA contingency volume~\cite{jarus_sora25} \\
\Qpop & a distribution of operations &
PAC-Bayes~\cite{mcallester1999pac}; regret~\cite{arora2012mwu}; mission
completion rate & per-flight-hour safety objectives; generalisation over the
ODD~\cite{easa_ai_concept2} \\
\bottomrule
\end{tabularx}
\end{center}

\begin{definition}[Sound transfer]
A guarantee at level $Q_g$ \emph{discharges} an objective at level $Q_o$ only if
$Q_g$ quantifies over a superset of what $Q_o$ requires and under a perturbation
model that covers the objective's hazard. Otherwise it is at most
\emph{evidence}, or it is a \emph{category error}.
\end{definition}

\begin{definition}[Lift]
A \emph{lift} is an argument that carries a guarantee from one level to another
--- \Qin$\to$\Qtr{} by a dynamics bound, \Qtr$\to$\Qop{} by a horizon and
re-anchoring argument, \Qop$\to$\Qpop{} by a concentration or PAC argument. Every
lift introduces an assumption, and that assumption is itself an assurance
obligation.
\end{definition}

\noindent
The parent manuscript's composition chain is a worked instance:
\Qin{} (a perturbation bound on the navigation input) lifted to \Qtr{} by
Gr\"onwall under a Lipschitz hypothesis, to \Qop{} by a corridor predicate under
re-anchoring, and to \Qpop{} by a mission distribution. Each of its three stated
conditions is exactly one lift's assumption. That is why it serves as the SoK's
running example rather than its subject.

% =====================================================================
\section{Hypotheses the systematisation tests}
% =====================================================================
These are predictions, and the paper is only worth writing if any of them can
come out false.

\begin{description}[leftmargin=2.4em,style=nextline]
\item[H1] \textbf{Most certified-robustness work for learned perception sits at
  \Qin, and therefore discharges no assurance objective without a lift.} The
  corpus will show the proportion.
\item[H2] \textbf{Lifts exist in the literature but are rarely accounted for as
  obligations.} Papers that lift a guarantee seldom list the lift's assumption as
  something requiring evidence. The coding scheme records this per paper.
\item[H3] \textbf{The \Qpop{} objectives --- per flight hour, generalisation
  over the ODD --- are the natural match for statistical guarantees, yet
  statistical guarantees for learned perception are the least developed
  class.} If true, this is the research agenda the SoK hands the community.
\item[H4] \textbf{Runtime assurance is the dominant practical bridge} for
  objectives no offline certificate discharges, because it converts a
  per-operation obligation into a per-input check plus a verified fallback.
\end{description}

\begin{honestly}
H1 could fail. If the corpus turns out to contain substantially more \Qtr{} and
\Qpop{} work than expected --- closed-loop reachability for neural controllers
is an active area --- the gap is narrower than this proposal claims, and the
paper must report that. An SoK whose finding is ``it mostly lines up'' would be
less exciting and still publishable, but only if it says so rather than choosing
a corpus that confirms the thesis.
\end{honestly}

% =====================================================================
\section{Method}
\label{sec:method}
% =====================================================================
\subsection{The guarantee corpus}
\begin{itemize}
\item \textbf{Sources.} IEEE S\&P, ACM CCS, USENIX Security, NDSS; NeurIPS, ICML,
  ICLR; CAV, HSCC, ICCPS, EMSOFT; DASC and the SAE/AIAA avionics venues. Years
  2017--2026.
\item \textbf{Inclusion.} A formal or statistical guarantee on a learned
  component used, or evaluated, for perception, state estimation, navigation or
  control of an autonomous vehicle. Benchmarks without a guarantee are out;
  guarantees on non-learned components are out.
\item \textbf{Search.} Keyword search seeded from the canonical systems
  (\cite{katz2017reluplex,wang2021betacrown,cohen2019smoothing,julian2016acasxu}),
  then forward and backward citation snowballing until saturation.
\item \textbf{Expected size.} Of the order of 80--150 papers. This is a
  prediction, not a count: \textbf{the corpus has not been collected.}
\end{itemize}

\subsection{The objective corpus}
SORA~2.5 operational safety objectives and the contingency-volume definition
in Annex~A~\cite{jarus_sora25}; the EASA AI Concept Paper Issue~2
learning-assurance objectives~\cite{easa_ai_concept2}; the MLEAP
findings~\cite{easa_mleap}; and ED-324 / ARP6983~\cite{ed324} at the level of
concepts only, since it is in preparation.

\subsection{Coding scheme}
Each guarantee and each objective is coded on the same fields, so they can be
joined:
\begin{center}\small
\begin{tabularx}{\linewidth}{@{}P{0.24\linewidth}Y@{}}
\toprule
\textbf{Field} & \textbf{Values} \\
\midrule
Quantifier level & \Qin, \Qtr, \Qop, \Qpop \\
Perturbation model & norm-bounded input; sensor delay; sensor bias; adversarial
spoofing; distribution shift; none \\
Lifts used & dynamics, horizon/re-anchoring, concentration, runtime monitor;
each with its stated assumption \\
Assumption accounted for? & yes / stated but not evidenced / unstated \\
Evidence class & proof; measured; simulated; modelled \\
Domain & perception, estimation, navigation, control, end-to-end \\
\bottomrule
\end{tabularx}
\end{center}

\noindent
Two coders code independently; disagreements are resolved by discussion, and
Cohen's $\kappa$ is reported for the quantifier and assumption fields, which are
the ones the findings rest on.

\subsection{The transfer table}
The output joins the two corpora on quantifier level and perturbation model and
classifies every (guarantee class, objective) pair as \textbf{discharge},
\textbf{evidence via lift} (naming the lift), or \textbf{category error}. P3's
discharge / evidence / untouched gap analysis becomes one slice of this table.

% =====================================================================
\section{The evidence generator}
% =====================================================================
Carried over from P3, and now with a sharper purpose: it demonstrates one
\emph{sound} transfer path end to end. From a certified run of the parent
project it emits a SORA Annex~A-shaped evidence pack in which
\begin{itemize}
\item the Operational Design Domain is the parent's \emph{measured} operating
  region (attitudes off singularity, airspeed below stall, actuators unsaturated),
  which P3 noticed is an ODD in all but name;
\item the contingency extent is the certified tube radius;
\item each lift's assumption appears as an explicit, open obligation rather than
  disappearing into the result; and
\item every figure carries its provenance tag.
\end{itemize}
The third point is the generator's real contribution. A compliance pack that
lists the certificate but hides its assumptions would be precisely the category
error the SoK exists to name.

% =====================================================================
\section{What the SoK delivers}
% =====================================================================
\begin{itemize}
\item \textbf{A quantifier framework} that places certified guarantees and
  assurance objectives on one axis.
\item \textbf{A coded corpus} of certified guarantees for learned autonomy,
  released with the coding scheme and agreement statistics.
\item \textbf{The transfer table}: which guarantees discharge which objectives,
  which need which lift, and which pairings are category errors.
\item \textbf{An accounting of lifts as obligations}, testing whether the
  literature treats them as such.
\item \textbf{A research agenda}: the empty cells of the transfer table,
  ranked by how many objectives they would discharge.
\item \textbf{An evidence generator} demonstrating one sound path.
\end{itemize}

% =====================================================================
\section{Relation to the parent manuscript}
% =====================================================================
P3 already said that it closes none of the parent's limitations, and the SoK
does not change that. What it changes is how one of them reads. The parent's
guarantee bounds a population of missions rather than a single flight, which a
reviewer can read as a weakness. The SoK shows that this is a property of every
\Qpop{} guarantee and of every per-flight-hour objective, so it is the correct
level for an airworthiness argument rather than a compromise. That is a
reframing, and it should be presented as one, not as a closure.

% =====================================================================
\section{Risks}
% =====================================================================
\begin{itemize}
\item \textbf{A thin or biased corpus.} The main way an SoK fails at S\&P. Stated
  inclusion criteria, snowballing to saturation, two coders and a reported
  $\kappa$ are the defence. If time forces a smaller corpus, the paper says
  how much smaller and drops the claims that need the full one.
\item \textbf{The standards move.} ED-324 is in preparation. Key every claim to a
  stable concept --- ODD, contingency volume, learning assurance --- and confine
  clause numbers to an appendix.
\item \textbf{Regulatory over-claim.} Nothing here certifies an aircraft. The
  paper maps mathematical guarantees onto an evidence framework, and the
  abstract must already say so.
\item \textbf{Being read as a survey.} The quantifier framework and the transfer
  table are what distinguish an SoK from a survey. If either is weak, the
  fallback is ACM Computing Surveys, and the paper should be retargeted rather
  than oversold.
\end{itemize}

% =====================================================================
\section{Page plan and timeline}
% =====================================================================
\begin{center}\small
\begin{tabular}{@{}lr@{\qquad}lr@{}}
\toprule
\textbf{Section} & \textbf{pp} & \textbf{Section} & \textbf{pp} \\
\midrule
Introduction                 & 1.25 & The transfer table          & 2.0 \\
The quantifier framework     & 1.75 & Lifts as obligations        & 1.25 \\
Method and corpus            & 1.5  & Evidence generator (case)   & 1.0 \\
Guarantees by level          & 1.75 & Research agenda             & 1.0 \\
Objectives by level          & 1.0  & Related SoKs, conclusion    & 0.5 \\
\midrule
\multicolumn{3}{@{}l}{\textbf{Body total}} & \textbf{13.0} \\
\bottomrule
\end{tabular}
\end{center}

\noindent
About three months, and it needs no measurement, so it can start at once and
run alongside Paper~I's hardware campaign. Weeks 1--6: search and coding.
Weeks 5--8: transfer table and generator. Weeks 9--12: writing. Submit it
before Paper~I, so that Paper~I can cite its vocabulary.

\small
\bibliographystyle{plain}
\bibliography{proposals}

\end{document}
